\scanpage{192}

et les

\[
\boxed{B_i\in\mathcal A_1^0=k\{X_0,X_1\}\qquad (i\geq0)}
\]

par les formules

\[
\left\{\begin{aligned}
\varphi(X_0^{(i)})&=(X_0+f)^{(i)}+B_i\,dX_0\\
\varphi(dX_0)&=1+df
\end{aligned}\right.
\]

(de sorte qu'on a

\[
\begin{aligned}
 f&=\varphi_{\deleted{\unsure{0}}}^{0}(X_0)-X_0=-\bigl(\varphi_{\deleted{\unsure{0}}}^{0}(X_1)-X_1\bigr)\\
 B_i\,d\deleted{\unsure{f}}X_0^{\unsure{\prime}}&=\varphi^1\bigl(X_0^{(i)}\bigr)
\end{aligned}
\]

\uline{NB}. Si on échange le rôle de $X_0$ et $X_1$, $f$ change de signe, - que deviennent les $B_i$ ? Il faut calculer

\[
\begin{aligned}
\varphi^1(X_{\deleted{\unsure{0}}1}^{(i)})&=\varphi^1\bigl((T-X_0)^{(i)}\bigr)\\
&=\varphi^1\biggl(\sum_{0\leq j\leq i}T^{(i-j)}X_0^{(j)}(-1)^j\biggr)
=\biggl(\sum_{0\leq j\leq i} T^{(i-j)}(-1)^j B_j\biggr)dX_0\\
&=\underbrace{\deleted{\unsure{-}}\biggl(\sum_{0\leq j\leq i}(-1)^{j+1} T^{(i-j)} B_j\biggr)}_{C_i}dX_1
\end{aligned}
\]

\deleted{\unsure{En termes|En tout cas}} \deleted{$\varphi^1(T)$}

\[
\begin{aligned}
&\underbrace{\varphi^1}_{\unsure{\pi_0^1\varphi|k'_0\varphi}}
\biggl(\underbrace{(X_0+X_1^{\unsure{\prime}})^{(i)}}_{\displaystyle\sum_{j+k=i}X_0^{(j)}X_1^{(k)}}\deleted{\unsure{=}}\biggr)
=\varphi^1(T^i)=0\\
&\underbrace{\hspace{4.4cm}}_{\displaystyle\pi^1\sum_{j+k=i}
\underbrace{\varphi(X_0^{(j)})}_{(X_0+f)^{(j)}+B_jdX_0}
\underbrace{\varphi(X_1^{(k)})}_{(X_1-f)^{(k)}+C_k\underbrace{dX_1}_{-dX_0}}}
\end{aligned}
\]

donc

\[
\sum_{j+k=i}\bigl[-(X_0+f)^{(j)}C_k\deleted{\unsure{-}}+B_j(X_1-f)^{(k)}\bigr]=0
\]

% archive sheet number: 191
